\documentclass[11pt,a4paper]{article}

\usepackage[margin=1in]{geometry}
\usepackage[T1]{fontenc}
\usepackage{lmodern}
\usepackage{microtype}
\usepackage{amsmath,amssymb,amsthm,mathtools}
\usepackage{booktabs,array,tabularx}
\usepackage{enumitem}
\usepackage{xcolor}
\usepackage{float}
\usepackage{hyperref}
\usepackage[nameinlink,capitalize,noabbrev]{cleveref}

\hypersetup{colorlinks=true,linkcolor=blue!45!black,
  citecolor=blue!45!black,urlcolor=blue!45!black,
  pdftitle={Near-Optimal Smooth Nonconvex Optimization with an Unknown Tail Index},
  pdfauthor={Colombo Team}}
\allowdisplaybreaks[2]
\setlength{\parskip}{0.25em}
\setlength{\parindent}{1.15em}
\setlist[itemize]{leftmargin=*,itemsep=0.25em,topsep=0.4em}
\setlist[enumerate]{leftmargin=*,itemsep=0.35em,topsep=0.4em}
\emergencystretch=2em
\renewcommand{\arraystretch}{1.2}
\floatstyle{ruled}
\newfloat{algorithm}{tbp}{loa}
\floatname{algorithm}{Algorithm}

\newcommand{\R}{\mathbb{R}}
\newcommand{\E}{\mathbb{E}}
\newcommand{\Prb}{\mathbb{P}}
\newcommand{\eps}{\varepsilon}
\newcommand{\cF}{\mathcal{F}}
\newcommand{\cP}{\mathcal{P}}
\newcommand{\cT}{\mathcal{T}}
\newcommand{\cG}{\mathcal{G}}
\newcommand{\cA}{\mathcal{A}}
\newcommand{\norm}[1]{\left\lVert #1\right\rVert}
\newcommand{\inner}[2]{\left\langle #1,#2\right\rangle}
\newcommand{\ind}{\mathbf{1}}
\newcommand{\FB}{\mathsf{BatchNSGD}}
\newcommand{\OtwoNC}{\mathsf{O2NC}}
\newcommand{\OAdaGradR}{\mathsf{OAdaGradR}}
\DeclareMathOperator*{\argmin}{arg\,min}
\DeclareMathOperator*{\argmax}{arg\,max}

\newtheorem{theorem}{Theorem}[section]
\newtheorem{proposition}[theorem]{Proposition}
\newtheorem{lemma}[theorem]{Lemma}
\newtheorem{corollary}[theorem]{Corollary}
\theoremstyle{definition}
\newtheorem{definition}[theorem]{Definition}
\newtheorem{assumption}[theorem]{Assumption}
\theoremstyle{remark}
\newtheorem{remark}[theorem]{Remark}
\crefname{proposition}{Proposition}{Propositions}
\crefname{lemma}{Lemma}{Lemmas}
\crefname{corollary}{Corollary}{Corollaries}
\crefname{definition}{Definition}{Definitions}
\crefname{assumption}{Assumption}{Assumptions}
\crefname{algorithm}{Algorithm}{Algorithms}

\title{Near-Optimal Smooth Nonconvex Optimization\\
with an Unknown Tail Index}
\author{Colombo Team\\[0.25em]
\small\href{mailto:colombo.2026@outlook.com}{\texttt{colombo.2026@outlook.com}}}
\date{November 23, 2026}

\begin{document}
\maketitle

\begin{abstract}
We study smooth nonconvex optimization with unbiased stochastic gradients whose
noise has a bounded centered $p$th moment for an unknown $p\in(1,2]$.
We construct a single algorithm that, given a query budget and bounds on
smoothness and the initial objective gap, uses neither $p$ nor the noise scale.
In the normalized model, its output satisfies
$\E\norm{\nabla F(\widehat x_T)}
\le C_p(\log(eT)/T)^{(p-1)/(3p-2)}$ for every fixed $p$.
The algorithm runs ordinary minibatch normalized gradient descent over a
dyadic grid of batch sizes and selects a candidate using fresh empirical
gradients. A martingale moment bound makes this selection rule valid without
tail estimation, including for conditionally unbiased, history-dependent noise.
All training and validation queries are included in the budget.
Combined with the existing lower bound for arbitrary randomized methods,
the result gives sample complexity
$O_p(\eps^{-(3p-2)/(p-1)}\log(e/\eps))$, within one logarithmic
factor of the known-index optimum. We also show that validation can use
a vanishing fraction of the total budget and establish a sharp logarithmic
barrier for a specified class of one-sided rate-grid certificates.
The latter is not a lower bound for general adaptive algorithms.
An exact rate without the logarithmic factor remains unresolved here.
\end{abstract}

\noindent\textbf{Keywords.} Stochastic nonconvex optimization; heavy-tailed noise;
unknown tail index; normalized gradient descent; oracle complexity; adaptation.

\begin{samepage}
\noindent\textbf{AI Usage.}
Nearly the entire research pipeline for this paper was carried out by
Colombo, powered by GPT-5.6 Sol.
\par
\end{samepage}

\section{Introduction}\label{sec:introduction}

The moment assumptions imposed on a stochastic gradient determine how much
averaging is needed before its direction becomes useful. When the noise has
only a finite $p$th moment, with $1<p<2$, a sample mean remains consistent,
but the familiar square-root averaging rate cannot in general be guaranteed. In smooth
nonconvex optimization this changes the optimal expected-gradient rate from
$T^{-1/4}$ to $T^{-(p-1)/(3p-2)}$. Normalization and
online-to-nonconvex conversion attain the latter expected-gradient rate,
with optimal tuning that generally depends on $p$
\cite{hubler2025normalization,liuzhou2025,liu2026oco}.

The exponent $p$ is often unavailable to the algorithm. Moreover, a noise law
may satisfy several moment conditions at once, so identifying a unique tail
index is neither necessary nor a well-defined prerequisite for optimization.
The appropriate question is whether one procedure, with the same inputs and
update rules, can obtain the correct rate on every moment class.
Liu and Zhou \cite{liuzhou2025} and H\"ubler, Fatkhullin, and He
\cite{hubler2025normalization} obtained unknown-index guarantees with rate
$T^{-(p-1)/(2p)}$. For $1<p<2$, this exponent is smaller than the
known-index exponent.

This paper addresses a focused version of the adaptation problem: the
algorithm knows a query budget $T$ and bounds $L$ and $\Delta$ on smoothness
and the initial objective gap, but knows neither $p$ nor the noise scale
$\sigma$. We prove that a dyadic grid of batch sizes, followed by a fully
charged validation step, achieves the known-index exponent up to a logarithm.
The proof uses standard normalized descent and a standard martingale moment
inequality. Its essential observation is that the batch size is the only
unknown scale that must be searched in this setting.

\subsection{Main result and mechanism}

Write
\begin{equation}\label{eq:exponents-intro}
 a(p)=\frac{p-1}{3p-2},\qquad
 b(p)=\frac{p-1}{p},\qquad
 r(p)=\frac{1}{a(p)}=\frac{3p-2}{p-1}.
\end{equation}
With $n$ training queries and a fixed batch size $\tau$, minibatch normalized
gradient descent gives the bound
\begin{equation}\label{eq:intro-batch}
 \E\norm{\nabla F(Y_{\tau,n})}
 \le C\sqrt{\frac{L\Delta\tau}{n}}
       +C_p\sigma\tau^{-b(p)}.
\end{equation}
The first term accounts for the number of optimization steps; the second
accounts for estimation error within a batch. Crucially, the stepsize
realizing \eqref{eq:intro-batch} uses $L,\Delta,n,\tau$, but neither
$p$ nor $\sigma$. This is the familiar minibatch NSGD tradeoff
\cite{hubler2025normalization}, for which we give a complete proof under
the conditional-moment oracle used here.

In the normalized model $L=\Delta=\sigma=1$, balancing the terms gives
$\tau\asymp n^{p/(3p-2)}$ and rate $n^{-a(p)}$.
A dyadic grid covers this scale to within a factor of two. Dividing the
training budget among $O(\log T)$ runs therefore supplies a candidate of
quality $O_p((\log T/T)^{a(p)})$.
Selecting that candidate is the remaining issue: neither objective values
nor exact gradients are available.

We select the smallest norm of a fresh empirical gradient. For $M$ candidates
and $m$ validation queries at each candidate, the selection error is at most
$C_p\sigma M^{1/p}m^{-b(p)}$. Because
\begin{equation}\label{eq:exponent-gap}
 b(p)-a(p)=\frac{2(p-1)^2}{p(3p-2)}>0,
\end{equation}
validation is asymptotically less expensive than optimization for each fixed
$p>1$. Neither the estimator nor the selector uses the exponent appearing
in its analysis.

The resulting statement has the quantifier order
\[
 \exists\,\mathsf A\quad\forall\,p\in(1,2]\quad
 \exists\,C_p<\infty\quad\forall\,T\ge1:
 \quad\mathfrak R_T^{\mathsf A}(p)
 \le C_p\left(\frac{\log(eT)}{T}\right)^{a(p)}.
\]
The algorithm $\mathsf A$ is fixed before $p$ is chosen. The constants
and the point at which validation becomes a lower-order term may depend
on $p$. We do not assert uniform asymptotics as $p\downarrow1$.

The paper makes three contributions within this scope. First, it gives a
self-contained, one-dimensional batch-size reduction with explicit query
accounting. Second, it proves that a single logarithm in sample complexity
suffices and that fresh validation can use $o(T)$ queries. Third, it
separates a sharp limitation of one-sided rate-grid certificates from the
unresolved question of exact adaptation by general algorithms.
Appendix~\ref{sec:beyond-log} records a sufficient descent condition for
a possible exact-rate method; it is not an additional unconditional
convergence result.

\subsection{Comparison with related work}\label{sec:related}

\paragraph{Clipping and normalization.}
Early heavy-tail analyses established convergence of clipped methods
\cite{zhang2020adaptive,nguyen2023clipped}.
H\"ubler et al.\ \cite{hubler2025normalization} proved optimal tuned
minibatch NSGD bounds and an algorithm-specific limitation for prescribed
single-run schedules. Liu and Zhou \cite{liuzhou2025} analyzed normalized
momentum, including relaxed smoothness and more general noise models.
Our comparison uses their classical smooth, additive-noise specialization.
Their unknown-index guarantees require fewer problem parameters than our
main algorithm, whereas our guarantee has the sharper exponent for
$1<p<2$. At $p=2$, those earlier guarantees already attain $T^{-1/4}$;
our uniform construction adds a logarithmic overhead at that endpoint.

\paragraph{Grid search and validation.}
The use of restarted runs and empirical-gradient selection has direct
precedent in the tuning-free method of Khaled and Jin
\cite{khaled2024tuning}. Their nonconvex result uses a stronger noise
condition, parameter hints, and a high-probability criterion, so it is not
a like-for-like rate comparison with the finite-moment expectation bounds
here. We use the same general selection principle with a moment estimate
valid for martingale differences. The moment inequality and normalized
descent are established tools, rather than new components of this paper.

\paragraph{A closely related online result.}
The LongTails project reports a parameter-free heavy-tail stationarity
result based on restarted normalized-momentum runs and fresh empirical
gradients \cite{longtails2026stationarity}. For fixed $p$ and normalized
positive parameters, its displayed budget is
$O_p(\eps^{-r(p)}\log^2(e/\eps))$.
Its algorithm does not take $L,\Delta,p,\sigma$, or the horizon as inputs.
Our result concerns the narrower adaptation problem with $L,\Delta,T$
available, has one logarithm in the corresponding budget, and permits
conditional moment bounds without an IID assumption.
Table~\ref{tab:comparison} keeps these distinctions explicit. In particular,
the general idea of near-optimal unknown-index optimization is already
present in that online result.

\begin{table}[t]
\centering
\small
\setlength{\tabcolsep}{4pt}
\begin{tabularx}{\textwidth}{@{}
>{\raggedright\arraybackslash}p{0.28\textwidth}
>{\raggedright\arraybackslash}p{0.20\textwidth}
>{\raggedright\arraybackslash}p{0.17\textwidth}
>{\raggedright\arraybackslash}X@{}}
\toprule
Result & Tuning information & Oracle / scope & Queries for
$\E\norm{\nabla F(\widehat x)}\le\eps$\\
\midrule
H\"ubler et al.\ \cite{hubler2025normalization}, Cor.~3;
Liu--Zhou \cite{liuzhou2025}, Thm.~3.2
& $p$ and problem parameters available
& IID (H\"ubler et al.); conditional (Liu--Zhou)
& $O_p(\eps^{-r(p)})$\\[0.35em]
Liu--Zhou \cite{liuzhou2025}, Thm.~3.6
& Horizon only in the classical specialization
& Conditional centered $p$th moment
& $O_p(\eps^{-2p/(p-1)})$\\[0.35em]
H\"ubler et al.\ \cite{hubler2025normalization}, Eq.~(4)
& Prescribed horizon; no $p,\sigma,L,\Delta$
& IID centered $p$th moment
& $O_p(\eps^{-2p/(p-1)})$\\[0.35em]
LongTails \cite{longtails2026stationarity}, online result
& No $p,\sigma,L,\Delta,T$
& IID centered $p$th moment
& $O_p(\eps^{-r(p)}\log^2(e/\eps))$\\[0.35em]
This paper, \cref{thm:main-upper,cor:sample-complexity}
& $L,\Delta,T$; no $p,\sigma$
& Conditional centered $p$th moment
& $O_p(\eps^{-r(p)}\log(e/\eps))$\\[0.35em]
Liu \cite{liu2026oco}, Thm.~11
& Even with $p$ known
& All randomized methods; IID hard oracle
& $\Omega_p(\eps^{-r(p)})$\\
\bottomrule
\end{tabularx}
\caption{Comparison on $L=\Delta=\sigma=1$, for each fixed
$p\in(1,2]$ and sufficiently small $\eps$. All costs count individual
stochastic-gradient calls and all rates concern the expected gradient
norm. The Liu--Zhou rows specialize their results to $L_1=\sigma_1=0$.
The LongTails row is the normalized asymptotic specialization of its
displayed online budget, consulted November 23, 2026; it is not a claim
that the general parameter dependence coincides. The lower bound permits
dimension to grow with accuracy.}\label{tab:comparison}
\end{table}

\paragraph{Online-to-nonconvex conversion and lower bounds.}
The conversion of Cutkosky, Mehta, and Orabona
\cite{cutkosky2023o2nc}, combined with Liu's heavy-tail optimistic
online-learning analysis \cite{liu2026oco}, provides another candidate
generator with the same batch-scale tradeoff. We retain that construction
in Appendix~\ref{app:o2nc}; the main proof uses minibatch NSGD to avoid
an unnecessary dependence on online regret machinery.
For optimality, we invoke Liu's lower bound for arbitrary randomized
methods \cite[Theorem~11]{liu2026oco}. Its algorithmic scope matters:
a lower bound restricted to zero-respecting methods would not, by itself,
justify an unrestricted minimax comparison. The finite-variance endpoint
is also consistent with the stochastic lower bounds of Arjevani et al.
\cite{arjevani2023lower}.

\section{Oracle model and performance criteria}\label{sec:model}

Let $F:\R^d\to\R$ be differentiable and lower bounded, with
$F_\star=\inf_xF(x)$. We are given $x_0\in\R^d$, a smoothness bound
$L>0$, and a gap bound $\Delta>0$ satisfying
\begin{equation}\label{eq:objective}
 F(x_0)-F_\star\le\Delta,\qquad
 \norm{\nabla F(x)-\nabla F(y)}\le L\norm{x-y}
 \quad(x,y\in\R^d).
\end{equation}
At call $t$, the query $x_t$ is measurable with respect to the history
$\cF_{t-1}$, which includes the algorithm's randomness available before
the response. The oracle returns $g_t$ with
\begin{equation}\label{eq:oracle}
 \E[g_t\mid\cF_{t-1}]=\nabla F(x_t),\qquad
 \E[\norm{g_t-\nabla F(x_t)}^p\mid\cF_{t-1}]\le\sigma^p
 \quad\text{a.s.}
\end{equation}
Here $p\in(1,2]$ and $\sigma\ge0$ are unknown to the algorithm.
The term \emph{fresh query} means a new response obeying
\eqref{eq:oracle} conditional on the complete current history; it does
not impose independence beyond this condition. IID stochastic gradients
are a special case.

The algorithm receives no function values. It cannot evaluate one random
sample at several points, and no smoothness of sample gradients is assumed.
Every call, including one made for output selection, counts toward the
budget. The output criterion is $\E\norm{\nabla F(\widehat x)}$.
The model allows infinite noise variance when $p<2$.

\begin{definition}[One common algorithm family]
A common family is a fixed collection
$\mathsf A=\{\mathsf A_T:T\in\mathbb N\}$ whose inputs may include
$T,L,\Delta,x_0$, but not $p$ or $\sigma$. Its formulas and its
rule for choosing queries are the same on every moment class.
\end{definition}

Let $\cP_p$ be the union, over all finite dimensions, of instances
satisfying \eqref{eq:objective}--\eqref{eq:oracle} with
$L=\Delta=\sigma=1$. For a common family, define
\begin{equation}\label{eq:risk-definition}
 \mathfrak R_T^{\mathsf A}(p)
 =\sup_{\mathcal I\in\cP_p}
 \E_{\mathcal I,\mathsf A}\norm{\nabla F(\widehat x_T)},\qquad
 N_p^{\mathsf A}(\eps)
 =\inf\{T\in\mathbb N:\mathfrak R_T^{\mathsf A}(p)\le\eps\}.
\end{equation}
Thus $N_p^{\mathsf A}(\eps)$ is a pointwise inversion of a
fixed-budget guarantee. It is not a stopping rule that chooses $T$
without $p$. Let $N_p^\star(\eps)$ denote the usual minimax query
complexity on $\cP_p$, allowing the algorithm to know $p$.
The upper bounds below hold in every dimension; their matching lower
bounds concern the union over dimensions.

We use $a(p),b(p),r(p)$ from \eqref{eq:exponents-intro} and
\begin{equation}\label{eq:identities}
 \alpha(p):=\frac{p}{3p-2}=1-2a(p),\qquad
 \alpha(p)b(p)=a(p),\qquad \frac1p+b(p)=1.
\end{equation}
Constants denoted by $C$ are universal; $C_p,c_p$ may depend only on
$p$. Their values may change between displays. All unmarked logarithms
are natural.

\begin{lemma}[Initial-gradient bound]\label[lemma]{lem:initial-gradient}
Under \eqref{eq:objective}, $\norm{\nabla F(x_0)}\le\sqrt{2L\Delta}$.
\end{lemma}
\begin{proof}
The descent lemma gives
$F_\star\le F(x_0-L^{-1}\nabla F(x_0))
\le F(x_0)-\norm{\nabla F(x_0)}^2/(2L)$.
\end{proof}
If a supplied gap bound is zero, this lemma allows immediate termination.
If $L=0$, the gradient is constant and lower boundedness forces it to be
zero, so this case is trivial as well.

\section{Validation without the tail index}\label{sec:validation}

The following estimate is a Hilbert-space martingale moment inequality.
We include a proof to make the conditional oracle argument explicit;
sharper constants are available, for example in
\cite[Lemma~10]{hubler2025normalization}.

\begin{lemma}[Empirical means of martingale differences]\label[lemma]{lem:martingale-mean}
Let $\xi_1,\ldots,\xi_m$ be Euclidean martingale differences
relative to $(\cF_i)_{i=0}^m$, with
$\E[\norm{\xi_i}^p\mid\cF_{i-1}]\le\sigma^p$ for
$1<p\le2$. There is a numerical constant $M_p\le16$ such that
\begin{equation}\label{eq:martingale-moment}
 \E\left[\left\|\frac1m\sum_{i=1}^m\xi_i\right\|^p
          \,\middle|\,\cF_0\right]
 \le M_p\sigma^p m^{1-p}.
\end{equation}
The constant does not depend on the dimension.
\end{lemma}
\begin{proof}
For $1<p\le2$,
\begin{equation}\label{eq:norm-smoothness}
 \norm{x+y}^p\le\norm{x}^p
 +p\norm{x}^{p-2}\inner{x}{y}+16\norm{y}^p,
\end{equation}
with the linear term set to zero at $x=0$.
For completeness, if $\norm y\ge\norm x/2$, the remainder after
subtracting the first two terms is at most
$(3^p+p2^{p-1})\norm y^p\le13\norm y^p$.
If $\norm y<\norm x/2$, the Hessian of $z\mapsto\norm z^p$
has operator norm at most $p\norm z^{p-2}$ on the segment
$x+[0,1]y$. Taylor's formula bounds the remainder by
$(p/2)(\norm x/2)^{p-2}\norm y^2\le\norm y^p$.
This proves \eqref{eq:norm-smoothness} in both cases.

Apply \eqref{eq:norm-smoothness} to the partial sums
$S_i=\sum_{j=1}^i\xi_j$. Conditional expectation removes the linear
term, and iteration gives
$\E[\norm{S_m}^p\mid\cF_0]\le16m\sigma^p$.
Division by $m^p$ proves the claim. The relevant products are integrable
by H\"older's inequality and induction on $i$.
\end{proof}

\begin{lemma}[Selection by fresh empirical gradients]\label[lemma]{lem:selection}
Let $Y_1,\ldots,Y_M$ be candidates determined before validation begins.
At each candidate obtain $m$ fresh gradients and set
\[
 \bar g_j=\frac1m\sum_{s=1}^m g_{j,s},\qquad
 \widehat j\in\argmin_{1\le j\le M}\norm{\bar g_j},
\]
breaking ties by the smallest index. Under \eqref{eq:oracle},
\begin{equation}\label{eq:selection}
 \E\norm{\nabla F(Y_{\widehat j})}
 \le \min_{1\le j\le M}\E\norm{\nabla F(Y_j)}
       +2M_p^{1/p}\sigma M^{1/p}m^{-b(p)}.
\end{equation}
The candidates and the validation errors need not be independent.
\end{lemma}
\begin{proof}
Put $e_j=\norm{\bar g_j-\nabla F(Y_j)}$. For any fixed index $j$,
\[
 \norm{\nabla F(Y_{\widehat j})}
 \le\norm{\bar g_{\widehat j}}+e_{\widehat j}
 \le\norm{\nabla F(Y_j)}+e_j+e_{\widehat j}
 \le\norm{\nabla F(Y_j)}+2\max_i e_i.
\]
During candidate $i$'s validation, its query point is fixed.
Conditioning on the history before that batch, applying
\cref{lem:martingale-mean}, and then taking expectations gives
$\E e_i^p\le M_p\sigma^pm^{1-p}$.
Consequently,
\[
 \E\max_i e_i
 \le\E\left(\sum_i e_i^p\right)^{1/p}
 \le\left(\sum_i\E e_i^p\right)^{1/p}
 \le M_p^{1/p}\sigma M^{1/p}m^{-b(p)}.
\]
Taking expectations in the pointwise comparison and minimizing over $j$
proves \eqref{eq:selection}.
\end{proof}

This argument does not need a good-candidate event or a bound on the
largest candidate gradient. It directly controls the expected loss from
selection, even when some candidates perform poorly.

\section{A fixed-batch candidate with no tail tuning}\label{sec:fixed-block}

We now give the elementary candidate generator used in the main algorithm.
It is ordinary minibatch normalized SGD, with its stepsize determined
by the number of updates. The algorithm and its tradeoff are standard
\cite{hubler2025normalization}; the proof below records precisely the
version required for our reduction.

\begin{algorithm}[H]
\caption{Fixed-batch normalized descent $\FB(x_0,L,\Delta;n,\tau)$}
\label{alg:batch}
\textbf{Input:} $L,\Delta>0$; integers $n\ge8$ and $1\le\tau\le n/2$.
\begin{enumerate}[label=\arabic*.,itemsep=0.2em,topsep=0.3em]
\item Set $K=\lfloor n/\tau\rfloor$,
$\eta=\sqrt{2\Delta/(LK)}$, and $x^{(0)}=x_0$.
\item For $k=0,\ldots,K-1$, obtain $\tau$ fresh gradients
$g_{k,1},\ldots,g_{k,\tau}$ at $x^{(k)}$. Set
\[
 v_k=\frac1\tau\sum_{s=1}^\tau g_{k,s},\qquad
 d_k=\begin{cases}v_k/\norm{v_k},&v_k\ne0,\\0,&v_k=0,\end{cases}
 \qquad x^{(k+1)}=x^{(k)}-\eta d_k.
\]
\item Draw $R$ uniformly from $\{0,\ldots,K-1\}$,
independently of the run, and return $Y_{\tau,n}=x^{(R)}$.
\end{enumerate}
\end{algorithm}

\begin{proposition}[Fixed-batch oracle inequality]\label[proposition]{prop:fixed-block}
Under \eqref{eq:objective}--\eqref{eq:oracle}, \cref{alg:batch} uses
$K\tau\le n$ queries, does not use $p$ or $\sigma$, and satisfies
\begin{align}
 \E\norm{\nabla F(Y_{\tau,n})}
 &\le\sqrt{\frac{2L\Delta}{K}}
       +2M_p^{1/p}\sigma\tau^{-b(p)}\notag\\
 &\le2\sqrt{\frac{L\Delta\tau}{n}}
       +2M_p^{1/p}\sigma\tau^{-b(p)}.
 \label{eq:fixed-block-bound}
\end{align}
\end{proposition}
\begin{proof}
For any vectors $h,v$, with $d=v/\norm v$ if $v\ne0$ and $d=0$
otherwise,
\begin{equation}\label{eq:direction}
 \inner{h}{d}\ge\norm h-2\norm{v-h}.
\end{equation}
For $v\ne0$, this follows from
$\inner{h}{d}\ge\norm v-\norm{v-h}$ and the triangle inequality;
the zero case is immediate. Let $h_k=\nabla F(x^{(k)})$.
By \cref{lem:martingale-mean} and conditional Jensen's inequality,
\[
 \E[\norm{v_k-h_k}\mid\text{history before batch }k]
 \le M_p^{1/p}\sigma\tau^{-b(p)}.
\]
Smoothness and \eqref{eq:direction} give
\[
 F(x^{(k+1)})\le F(x^{(k)})-\eta\norm{h_k}
       +2\eta\norm{v_k-h_k}+\frac{L\eta^2}{2}.
\]
Sum over $k$, take expectations, and use the initial gap to obtain
\[
 \frac1K\sum_{k=0}^{K-1}\E\norm{h_k}
 \le\frac{\Delta}{K\eta}+\frac{L\eta}{2}
       +2M_p^{1/p}\sigma\tau^{-b(p)}.
\]
Our choice of $\eta$ makes the first two terms equal and their sum
$\sqrt{2L\Delta/K}$. The random output realizes the average on the
left. Finally, $\tau\le n/2$ implies $K\ge n/(2\tau)$.
All quantities are integrable: the path length is at most $K\eta$,
and smoothness bounds the true gradients on that deterministic-radius ball.
\end{proof}

For a known $p$ and $\sigma>0$, the balancing scale is
\begin{equation}\label{eq:oracle-block}
 \tau_p(n)=\left(\frac{\sigma^2n}{L\Delta}\right)^{\alpha(p)}.
\end{equation}
When this scale lies within the feasible range, both terms of
\eqref{eq:fixed-block-bound} are of order
\begin{equation}\label{eq:balanced-value}
 \left(\frac{L\Delta\,\sigma^{p/(p-1)}}{n}\right)^{a(p)}.
\end{equation}
Equation~\eqref{eq:oracle-block} is used only in the analysis; no part
of the proposed algorithm evaluates it.

\section{Dyadic adaptation and its guarantees}\label{sec:main-upper}

\begin{algorithm}[H]
\caption{Dyadic batch-size portfolio $\mathsf A_{\rm dyad}$}
\label{alg:portfolio}
\textbf{Input:} Query budget $T\ge1$, point $x_0$, bounds $L,\Delta>0$.
\begin{enumerate}[label=\arabic*.,itemsep=0.25em,topsep=0.3em]
\item Set $J=\lceil\log_2(eT)\rceil$ and
$n=m=\lfloor T/(2J)\rfloor$. If $n<8$, return $x_0$.
\item Form $\cT_n=\{2^j:j\in\mathbb Z_{\ge0},\ 2^j\le n/2\}$.
For each $\tau\in\cT_n$, in increasing order, run
$Y_\tau=\FB(x_0,L,\Delta;n,\tau)$ using fresh queries.
\item After all training runs, obtain $m$ fresh gradients at each
$Y_\tau$ and form their mean $\bar g_\tau$.
\item Choose $\widehat\tau\in\argmin_{\tau\in\cT_n}\norm{\bar g_\tau}$,
breaking ties by the smaller batch size, and return
$\widehat x_T=Y_{\widehat\tau}$.
\end{enumerate}
\end{algorithm}

There are $M:=|\cT_n|\le J$ candidates. The total number of calls is
at most $M(n+m)\le T$, including all validation calls.
Each run starts from $x_0$; the runs need not be probabilistically
independent under \eqref{eq:oracle}.

\begin{theorem}[Near-optimal adaptation]\label{thm:main-upper}
In the normalized model, \cref{alg:portfolio} is one common family
such that, for every fixed $p\in(1,2]$,
\begin{equation}\label{eq:main-upper}
 \mathfrak R_T^{\mathsf A_{\rm dyad}}(p)
 \le C_p\left(\frac{\log(eT)}{T}\right)^{a(p)}
 \qquad(T\ge1).
\end{equation}
More explicitly, whenever $n\ge8$, its risk is at most
\begin{equation}\label{eq:finite-horizon}
 C\left(\frac JT\right)^{a(p)}+CJT^{-b(p)},
\end{equation}
where $C$ can be chosen independently of $p\in(1,2]$.
Every candidate generator and the selection rule use neither $p$ nor
$\sigma$.
\end{theorem}
\begin{proof}
Fix $p$, and abbreviate $a=a(p)$, $b=b(p)$, and $\alpha=\alpha(p)$.
Suppose first that $n\ge8$. Put
$u=\min\{n^\alpha,n/2\}$ and choose the largest dyadic
$\tau\le u$, so that $u/2<\tau\le u$.
If $n^\alpha\le n/2$, then
\[
 \sqrt{\tau/n}\le n^{-a},\qquad
 \tau^{-b}\le2^b n^{-\alpha b}=2^b n^{-a}.
\]
If $n^\alpha>n/2$, then $n^{-a}>1/\sqrt2$ and
$n/4<\tau\le n/2$. Consequently,
\[
 \sqrt{\tau/n}\le1/\sqrt2<n^{-a},\qquad
 \tau^{-b}\le4^b n^{-b}\le2n^{-a}.
\]
Here $0<b\le1/2$ and $b\ge a$.
In either case, \cref{prop:fixed-block} supplies a candidate with
\begin{equation}\label{eq:good-candidate}
 \E\norm{\nabla F(Y_\tau)}\le Cn^{-a},
\end{equation}
with a universal constant because $M_p^{1/p}\le16$.

By \cref{lem:selection}, the additional validation error is at most
$CM^{1/p}m^{-b}\le CJ^{1/p}n^{-b}$.
Since $n\ge8$ and $n=\lfloor T/(2J)\rfloor$, we have
$n\ge T/(4J)$. The identity $1/p+b=1$ now yields
\[
 \mathfrak R_T^{\mathsf A_{\rm dyad}}(p)
 \le C(J/T)^a+CJT^{-b},
\]
proving \eqref{eq:finite-horizon}.
For each fixed $p>1$, \eqref{eq:exponent-gap} implies
\[
 \frac{JT^{-b}}{(J/T)^a}
 =J^{1-a}T^{-(b-a)}\longrightarrow0.
\]
As $J\asymp\log(eT)$, this proves \eqref{eq:main-upper} for
$n\ge8$ after increasing $C_p$ to cover a finite initial range.
When $n<8$, \cref{lem:initial-gradient} bounds the returned gradient
by $\sqrt2$, which is absorbed in the same way.
\end{proof}

The two-term form \eqref{eq:finite-horizon} distinguishes the training
cost from the validation cost. In particular, it avoids hiding a
$p$-dependent crossover inside the algorithm itself.

\subsection{General parameters}\label{sec:general-parameters}

\begin{theorem}[Portfolio oracle inequality]\label{thm:parameter-oracle}
For each $\tau\in\cT_n$, generate $Y_{\tau,n}$ by
\cref{alg:batch} with at most $n\ge8$ training calls, and give every
candidate $m\ge1$ fresh validation calls.
For $M=|\cT_n|$, the resulting output satisfies
\begin{align}
 \E\norm{\nabla F(\widehat x)}
 \le {}&\min_{\tau\in\cT_n}
 \left\{2\sqrt{\frac{L\Delta\tau}{n}}
       +2M_p^{1/p}\sigma\tau^{-b(p)}\right\}\notag\\
 &+2M_p^{1/p}\sigma M^{1/p}m^{-b(p)}.
 \label{eq:parameter-oracle}
\end{align}
\end{theorem}
\begin{proof}
Apply \cref{lem:selection} to the candidate with the smallest bound
in \cref{prop:fixed-block}.
\end{proof}

If $2\le\tau_p(n)\le n/2$, dyadic rounding gives
\begin{equation}\label{eq:general-interior}
 \E\norm{\nabla F(\widehat x)}
 \le C_p\left[
 \left(\frac{L\Delta\,\sigma^{p/(p-1)}}{n}\right)^{a(p)}
 +\sigma M^{1/p}m^{-b(p)}\right].
\end{equation}
For fixed $L,\Delta,\sigma>0$ this interior condition eventually holds.
Outside it, \eqref{eq:parameter-oracle} remains valid, with a boundary
batch size available. If $\sigma=0$, choosing $\tau=1$ gives
$O(\sqrt{L\Delta/n})$ and validation is exact.
We do not turn the interior simplification into a uniform matching claim
over all joint regimes of $L,\Delta,\sigma,T$.

\subsection{Lower bound and sample complexity}\label{sec:lower}

We use the following existing lower bound, rather than a new hard-instance
construction. Liu \cite[Theorem~11]{liu2026oco} proves that, for
fixed $p\in(1,2]$, any randomized stochastic first-order method needs
\begin{equation}\label{eq:liu-lower}
 \Omega_p\left(L\Delta\eps^{-2}
       +L\Delta\sigma^{p/(p-1)}\eps^{-r(p)}\right)
\end{equation}
queries in the worst case to guarantee an expected gradient norm at most
$\eps$, for $\eps$ at most a universal constant times
$\sqrt{L\Delta}$. The oracle in the hard family satisfies the centered
moment condition, and the construction allows the dimension to grow.

\begin{corollary}[Pointwise near-matching complexity]\label[corollary]{cor:sample-complexity}
For every fixed $p\in(1,2]$ and sufficiently small $\eps>0$,
\begin{equation}\label{eq:sample-complexity}
 c_p\eps^{-r(p)}
 \le N_p^\star(\eps)
 \le N_p^{\mathsf A_{\rm dyad}}(\eps)
 \le C_p\eps^{-r(p)}\log\frac e\eps.
\end{equation}
The lower bound also has the following fixed-budget form: for every
sufficiently large $T$ and every randomized
$T$-query method, there is a normalized instance, in a dimension that
may depend on $T$, with expected output gradient at least
$c_pT^{-a(p)}$.
\end{corollary}
\begin{proof}
The first inequality is \eqref{eq:liu-lower} in the normalized model.
The middle inequality follows because the fixed-$p$ minimax class allows
the common algorithm. For the last inequality, substitute
$T=\lceil C_p\eps^{-r(p)}\log(e/\eps)\rceil$ into
\eqref{eq:main-upper} and choose the constant sufficiently large.
The fixed-budget lower statement follows by taking
$\eps=c_pT^{-a(p)}$ in \eqref{eq:liu-lower}.
\end{proof}

For example, at $p=3/2$ the earlier unknown-index exponent gives
$O(\eps^{-6})$, the present bound is
$O(\eps^{-5}\log(e/\eps))$, and the lower bound is
$\Omega(\eps^{-5})$. The polynomial gain is strict only for $p<2$,
since
\[
 \frac{2p}{p-1}-\frac{3p-2}{p-1}=\frac{2-p}{p-1}.
\]
No fixed-dimensional optimality statement is inferred from
\cref{cor:sample-complexity}.

\section{The role of the logarithm}\label{sec:logarithm}

\subsection{Validation can use a vanishing budget share}\label{sec:sublinear-validation}

The equal split between training and validation in \cref{alg:portfolio}
keeps its query accounting simple. It is not needed for the asymptotic
rate.

\begin{proposition}[A common sublinear validation schedule]\label[proposition]{prop:sublinear-validation}
Suppose $1\le M=O(\log T)$ candidates are available, and set
\begin{equation}\label{eq:sublinear-schedule}
 m_T=\left\lceil T\exp\{-\sqrt{\log(eT)}\}\right\rceil.
\end{equation}
Then $Mm_T=o(T)$ and, for every fixed $p>1$,
\begin{equation}\label{eq:sublinear-error}
 M^{1/p}m_T^{-b(p)}
 =o_p\bigl((M/T)^{a(p)}\bigr).
\end{equation}
Neither the schedule nor the selection rule uses $p$.
\end{proposition}
\begin{proof}
The total cost is at most
$MT\exp\{-\sqrt{\log(eT)}\}+M=o(T)$.
The ratio of the left side of \eqref{eq:sublinear-error} to its
comparison rate is bounded by
\[
 M^{1/p-a(p)}T^{-(b(p)-a(p))}
 \exp\{b(p)\sqrt{\log(eT)}\},
\]
which tends to zero since $b(p)>a(p)$ and $M=O(\log T)$.
\end{proof}

A concrete budgeted variant fixes
$n_0=\lfloor T/(4J)\rfloor$, uses the dyadic grid $\cT_{n_0}$,
and spends $m_T$ queries per candidate on validation. If
$n_0<8$ or $M(n_0+m_T)>T$, it returns $x_0$.
These fallback conditions cease to occur for sufficiently large $T$.
The grid still contains a constant-factor approximation to the optimal
scale for a training budget of order $T/\log T$, so the same main rate
follows. Thus validation does not require a positive asymptotic share
of the budget. In this construction, the leading logarithm arises from
training separate candidates at logarithmically many scales.

\subsection{What the certificate barrier establishes}

There is a precise sense in which a one-sided rate-grid proof cannot
remove the logarithm. Suppose a portfolio runs algorithms tuned to
guessed exponents and uses only guarantees of the form $n_j^{-s_j}$,
valid when $s_j\le s$ for the true rate coordinate $s$.
If the training budgets are disjoint and sum to $T$, then on any
nondegenerate interval $[u,v]\subset(0,1/4]$ the worst relative
certificate overhead is at least
$(\log T)^{v-o(1)}$. A uniform grid achieves
$O((\log T)^v)$. Appendix~\ref{sec:static-barrier} gives the exact
model and proof.

This statement concerns those certificates, not the realized risk of
every portfolio. In particular, the fixed-batch bound
\eqref{eq:fixed-block-bound} holds at every $p$ and has two terms; it
is not itself a one-sided guessed-exponent certificate. The barrier
therefore does not establish that the logarithm in our algorithm, or
in the unrestricted adaptation problem, is necessary.

\section{Discussion}\label{sec:discussion}

A finite grid suffices to recover the optimal heavy-tail exponent
without estimating the tail index. The reason is specific and simple:
the training tradeoff depends on one batching scale, while empirical
gradient validation converges faster than the target optimization rate
for every fixed $p>1$. The method uses elementary normalized descent,
and its upper bound is independent of dimension.

The main unresolved issue is whether a single common algorithm can
attain $O_p(T^{-a(p)})$ without a logarithmic loss. The static-certificate
result does not answer this question. Appendix~\ref{sec:beyond-log}
gives a shared-trajectory descent reduction in which exact adaptation
would follow from controlling two explicit error budgets. Those bounds
are sufficient conditions, not a characterization of all successful
algorithms. Likewise, small observed variation does not by itself
certify the accuracy of a heavy-tailed mean; the appendix records this
limitation without turning it into an impossibility theorem.

The present results concern expected stationarity and a prescribed
query budget. A high-probability guarantee, a stopping rule that
recognizes a requested accuracy without knowing $p$, and comparable
adaptation when $L$ and $\Delta$ are unavailable require separate
analysis. The stronger parameter adaptivity reported by LongTails
\cite{longtails2026stationarity} illustrates why these distinctions
should be kept separate from the precision exponent alone.

\appendix


\section{An alternative construction through online-to-nonconvex conversion}
\label{app:o2nc}

The main algorithm uses normalized minibatch descent.  An alternative
tail-index-free candidate follows from the online-to-nonconvex framework of
Cutkosky, Mehta, and Orabona~\cite{cutkosky2023o2nc} and Liu's optimistic
heavy-tail analysis~\cite{liu2026oco}.  We record the construction to make
its relationship with the fixed-block tradeoff explicit.  This appendix
assumes $L\Delta>0$, $n\ge8$, and $2\le\tau\le n/2$.  For $\tau=1$, use
the normalized minibatch candidate from \cref{sec:fixed-block} instead.

\paragraph{The online learner and the queried path.}
Let $K=\lfloor n/\tau\rfloor$, choose $\delta>0$, and set $D=\delta/\tau$.
Write $B_D=\{u\in\R^d:\norm u\le D\}$, and define its linear minimizer by
\[
 \mathsf{LM}_D(v)=
 \begin{cases}
  -Dv/\norm v,&v\ne0,\\
  0,&v=0.
 \end{cases}
\]
Initialize $y_0=x_0$, $u_1=0$, $g_0=0$, and $V_0=0$.
At round $s=1,\ldots,K\tau$, draw a fresh $q_s\sim\mathrm{Unif}[0,1]$,
query a fresh gradient $g_s$ at
\[
 z_s=y_{s-1}+q_su_s,
 \qquad y_s=y_{s-1}+u_s,
\]
and update
\[
 V_s=
 \begin{cases}
  \norm{g_s-g_{s-1}}^2,&\text{$s$ is the first round of a block},\\
  V_{s-1}+\norm{g_s-g_{s-1}}^2,&\text{otherwise}.
 \end{cases}
\]
If $s$ is the last round of a block, set $u_{s+1}=\mathsf{LM}_D(g_s)$.
Otherwise, set
\begin{equation}\label{eq:o2nc-online-update}
 u_{s+1}\in\argmin_{u\in B_D}
 \left\{\inner{2g_s-g_{s-1}}{u}
       +\frac{\sqrt{V_s}}{2\sqrt2D}\norm{u-u_s}^2\right\}.
\end{equation}
When $V_s>0$, this is the Euclidean projection of
$u_s-\sqrt2D(2g_s-g_{s-1})/\sqrt{V_s}$ onto $B_D$.
When $V_s=0$, use the well-defined linear minimizer
$\mathsf{LM}_D(2g_s-g_{s-1})$.

These are Liu's $\OAdaGradR$ updates (Algorithm~6), used inside $\OtwoNC$
(Algorithm~4).  The previous-gradient hint is retained across blocks;
only the variation accumulator is cleared.  The boundary action
$\mathsf{LM}_D(g_s)$ is essential to the stated regret guarantee.
Neither update uses $p$ or $\sigma$.
Define the block averages
\[
 \bar z_k=\frac1\tau\sum_{s=(k-1)\tau+1}^{k\tau}z_s,
 \qquad k=1,\ldots,K,
\]
and return one of them uniformly at random.  Exactly $K\tau\le n$
oracle calls are used.

\begin{proposition}[Liu's smooth specialization]\label[proposition]{prop:liu-specialization}
Under the oracle assumptions of this paper, the preceding construction satisfies
\begin{equation}\label{eq:liu-specialization}
 \E\!\left[\frac1K\sum_{k=1}^K\norm{\nabla F(\bar z_k)}\right]
 \le C_p\left[
  \frac{\norm{\nabla F(x_0)}}{K\tau}
  +\frac{\Delta}{\delta K}
  +\sigma\tau^{-b(p)}+L\delta\right].
\end{equation}
\end{proposition}

\begin{proof}
This is Theorem~17 of the full version of~\cite{liu2026oco}, with its
H\"older parameters set to $\nu=1$, $H=L$, and $G=0$, and its block
length $T$ identified with $\tau$.  The conditional moment assumption
here implies the moment assumption used there.
\end{proof}

\paragraph{Recovering the fixed-block bound.}
Choose $\delta=\sqrt{\Delta\tau/(Ln)}$.  Since $\tau\le n/2$,
$K\ge n/(2\tau)$ and $K\tau\ge n/2$.  Consequently,
\[
 \frac{\Delta}{\delta K}
 \le2\sqrt{\frac{L\Delta\tau}{n}},
 \qquad
 L\delta=\sqrt{\frac{L\Delta\tau}{n}}.
\]
For the returned point $Y^{\rm O2NC}_{\tau,n}$,
\begin{equation}\label{eq:o2nc-fixed-block-bound}
 \E\norm{\nabla F(Y^{\rm O2NC}_{\tau,n})}
 \le C_p\left[
  \frac{\norm{\nabla F(x_0)}}n
  +\sqrt{\frac{L\Delta\tau}{n}}
  +\sigma\tau^{-b(p)}\right].
\end{equation}
The initial-gradient term is at most $\sqrt{2L\Delta}/n$ and can be
absorbed into the square-root term.  Hence this alternative candidate
has the same tradeoff, and the same dyadic selection argument applies.
The normalized minibatch construction is used in the main text because
its proof and implementation are shorter.


\section{A sharp limitation of static rate-grid certificates}\label{sec:static-barrier}

This appendix concerns a specified class of numerical performance certificates.
It does not impose a restriction on arbitrary stochastic first-order algorithms.
In particular, the one-sided validity condition below describes a portfolio of
guessed tail exponents; it need not describe every block-length portfolio.

Write
\[
 s=a(q)=\frac{q-1}{3q-2}\in(0,1/4],
 \qquad q=\frac{1-2s}{1-3s},
\]
and fix a nondegenerate interval $I=[u,v]\subset(0,1/4]$.
At horizon $T$, a static rate-grid certificate consists of finitely many
coordinates $s_{T,j}\in(0,1/4]$ and training budgets $n_{T,j}\ge1$ satisfying
$\sum_j n_{T,j}\le T$. Candidate $j$ has certificate $n_{T,j}^{-s_{T,j}}$,
which is valid at true coordinate $s$ only when $s_{T,j}\le s$.
We study the resulting normalized overhead
\begin{equation}\label{eq:overhead}
 h_T(s):=\min_{j:s_{T,j}\le s}
 T^{s-s_{T,j}}\left(\frac{T}{n_{T,j}}\right)^{s_{T,j}},
 \qquad H_T(I):=\sup_{s\in I}h_T(s),
\end{equation}
where a minimum over the empty set is $+\infty$.
A common positive prefactor in all certificates changes $H_T(I)$ by that
prefactor and has no effect on the logarithmic exponent below.

\begin{lemma}[Weighted covering inequality]\label[lemma]{lem:weighted-cover}
For $r>0$, $R\ge1$, and any finite sequence $w_j>0$ with
$\sum_jw_j\le1$,
\[
 \sum_j[\log R+r\log w_j]_+\le\frac re R^{1/r}.
\]
\end{lemma}
\begin{proof}
Let $A=\{j:\log R+r\log w_j>0\}$ and $M=|A|$.
The claim is immediate if $M=0$. Otherwise, concavity of the logarithm gives
\[
 \sum_{j\in A}\log w_j
 \le M\log\left(\frac{\sum_{j\in A}w_j}{M}\right)
 \le-M\log M.
\]
The left side of the claimed inequality is at most
$M\log R-rM\log M$. Its maximum over real $M>0$ is
$(r/e)R^{1/r}$, attained at $M=R^{1/r}/e$.
\end{proof}

\begin{theorem}[Static certificate barrier]\label{thm:static-barrier}
For every sequence of the static certificates defined above,
\[
 \liminf_{T\to\infty}\frac{\log H_T(I)}{\log\log T}\ge v.
\]
Thus, on $I=[a(p_0),1/4]$ with $1<p_0<2$, the worst-case certificate
overhead is at least $(\log T)^{1/4-o(1)}$.
\end{theorem}
\begin{proof}
Set $L_T=\log T$, $h=(\log L_T)^{-1/2}$, and $R=H_T(I)$.
For sufficiently large $T$, the interval $J=[v-h,v]$ lies in $I$ and
$r:=v-3h/2$ is positive. Each finite overhead in \eqref{eq:overhead}
is at least one. We may assume $1\le R<\infty$.
Put $w_j=n_{T,j}/T$, so $\sum_jw_j\le1$.

Candidate $j$ covers a coordinate $s\ge s_{T,j}$ at overhead at most $R$
only if
\[
 (s-s_{T,j})L_T-s_{T,j}\log w_j\le\log R.
\]
Its covered interval therefore has length at most
$[\log R+s_{T,j}\log w_j]_+/L_T$.
If $R\ge T^{h/2}$, then
$\log R/\log\log T\ge hL_T/(2\log L_T)\to\infty$.
For all other horizons, let $B$ be the indices of candidates covering
at least one point of $J$. Since $\log w_j\le0$, such a candidate satisfies
$s-s_{T,j}\le\log R/L_T<h/2$ and hence $s_{T,j}\ge r$.
Their intervals cover $J$, so
\[
 hL_T
 \le\sum_{j\in B}[\log R+s_{T,j}\log w_j]_+
 \le\sum_j[\log R+r\log w_j]_+
 \le\frac reR^{1/r}.
\]
The last inequality is \cref{lem:weighted-cover}.
Consequently $R\ge(ehL_T/r)^r$.
As $r\to v$ and $\log h=o(\log L_T)$, this implies
$\log R\ge(v-o(1))\log\log T$. Together with the other case, this proves
the asserted lower limit.
\end{proof}

\begin{proposition}[Matching envelope]\label[proposition]{prop:static-upper}
For every fixed $I=[u,v]\subset(0,1/4]$, a grid of mesh $1/\log T$
with equal training budgets satisfies
$H_T(I)=O((\log T)^v)$.
\end{proposition}
\begin{proof}
Take the grid points $u+k/\log T$ that lie in $I$, together with $v$.
Their number is $K=O(\log T)$. Allocate $n_j=\lfloor T/K\rfloor$.
For all sufficiently large $T$, $n_j\ge T/(2K)$.
For each $s\in I$, there is a grid point $s_j\le s$ with
$s-s_j\le1/\log T$. Therefore
\[
 T^{s-s_j}(T/n_j)^{s_j}\le e(2K)^v=O((\log T)^v).
\]
\end{proof}

The two results characterize the worst-case logarithmic exponent of this
certificate model. They leave open the possibility that a different analysis,
sample sharing, or a common trajectory removes the logarithm in optimization.

\section{Shared-trajectory benchmarks and sufficient conditions}\label{sec:beyond-log}

The following results give a known-tail benchmark and a sufficient condition
for exact adaptation. They do not improve the unconditional guarantee of
\cref{thm:main-upper}. Their role is to state precisely what an adaptive
selection argument would still need to establish.

\subsection{Lagged moving-average normalized descent}

Let $T\ge2$, $1\le w\le T-1$, and $\eta>0$.
Set $x_1=x_0$ and $x_2=x_1$, and obtain fresh oracle responses $g_t$ at
$x_t$. For $t=2,\ldots,T$, define
\[
 \ell_t=\min\{w,t-1\},\qquad
 m_t=\frac1{\ell_t}\sum_{i=t-\ell_t}^{t-1}g_i,\qquad
 d_t=\begin{cases}m_t/\norm{m_t},&m_t\ne0,\\0,&m_t=0,\end{cases}
\]
and update $x_{t+1}=x_t-\eta d_t$.
The direction uses only past responses. The construction needs at most $T$
queries; the response at $x_T$ may be omitted because it is never used.

\begin{proposition}[Moving-average guarantee]\label[proposition]{prop:moving-average}
Under \eqref{eq:oracle},
\begin{equation}\label{eq:moving-average-bound}
 \frac1{T-1}\sum_{t=2}^T\E\norm{\nabla F(x_t)}
 \le\frac{\Delta}{\eta(T-1)}+L\eta(w+2)
 +C_p\sigma\left(w^{-b(p)}+\frac{w^{1/p}}T\right).
\end{equation}
\end{proposition}
\begin{proof}
Write $\xi_i=g_i-\nabla F(x_i)$. The martingale mean estimate
(\cref{lem:martingale-mean}) and Jensen's inequality give
\[
 \E\norm{\frac1{\ell_t}\sum_{i=t-\ell_t}^{t-1}\xi_i}
 \le C_p\sigma\ell_t^{-b(p)}.
\]
Since each displacement has length at most $\eta$,
\[
 \frac1{\ell_t}\sum_{i=t-\ell_t}^{t-1}
 \norm{\nabla F(x_i)-\nabla F(x_t)}
 \le\frac{L\eta}{\ell_t}\sum_{r=1}^{\ell_t}r
 =\frac{L\eta(\ell_t+1)}2.
\]
For $b=b(p)\in(0,1)$,
$\sum_{r=1}^w r^{-b}\le w^{1-b}/(1-b)$.
Splitting the initial windows from the full windows therefore yields
\[
 \frac1{T-1}\sum_{t=2}^T\ell_t^{-b(p)}
 \le C_p\left(w^{-b(p)}+\frac{w^{1/p}}T\right).
\]
For any vectors $h,m$, with $d=m/\norm m$ and the zero convention,
$\inner{h}{d}\ge\norm h-2\norm{h-m}$.
The descent lemma consequently gives
\[
 F(x_{t+1})\le F(x_t)-\eta\norm{\nabla F(x_t)}
 +2\eta\norm{m_t-\nabla F(x_t)}+\frac{L\eta^2}2.
\]
Sum from $t=2$ to $T$, use $x_2=x_0$ and the gap bound, and substitute
the preceding estimates. Since $\ell_t\le w$, the deterministic error
contributes at most $L\eta(w+3/2)\le L\eta(w+2)$ after division by
$\eta(T-1)$.
\end{proof}

\begin{corollary}[Known-tail shared-trajectory rate]\label[corollary]{cor:moving-average-rate}
In the normalized model, fix $p\in(1,2]$, put $s=a(p)$, and choose
\[
 w=\min\{T-1,\lceil T^{1-2s}\rceil\},\qquad \eta=T^{s-1}.
\]
An iterate sampled uniformly from $\{x_2,\ldots,x_T\}$ satisfies
$\E\norm{\nabla F(\widehat x_T)}\le C_pT^{-s}$.
\end{corollary}
\begin{proof}
For $T\ge2$, the chosen window lies between
$T^{1-2s}/2$ and $2T^{1-2s}$.
The identity $b(p)(1-2a(p))=a(p)$ gives
$w^{-b(p)}\le C_pT^{-s}$, while
$1/(\eta(T-1))\le2T^{-s}$ and $\eta(w+2)\le4T^{-s}$.
Finally, $w^{1/p}/T\le C_pT^{-3s}$.
Apply \cref{prop:moving-average} and average over the output index.
\end{proof}

\subsection{A thresholded selection criterion}

Fix $0<s_1<\cdots<s_K\le1/4$ and set
\begin{equation}\label{eq:rate-scales}
 c_j=T^{-s_j},\qquad \eta_j=T^{s_j-1}=\frac1{Tc_j},
 \qquad w_j=\lceil Tc_j^2\rceil.
\end{equation}
At round $t$, a predictable selector chooses $J_t\in\{0,1,\ldots,K\}$
and a predictor $m_{t,J_t}$ from the common history.
For $J_t\ge1$, use its normalized direction and update
$x_{t+1}=x_t-\eta_{J_t}d_t$.
The action $J_t=0$ means no update; set $\eta_0=0$ and $m_{t,0}=d_t=0$.
Start at $x_1=x_0$ and write
\[
 h_t=\nabla F(x_t),\qquad H_t=\norm{h_t},\qquad
 e_t=\norm{m_{t,J_t}-h_t}.
\]
For a comparator index $k$, define
\begin{align}
 \mathcal A_T&=\E\sum_{t=1}^T
 \eta_{J_t}(e_t-H_t/4)_+,\label{eq:A-budget}\\
 \mathcal Q_T&=\E\sum_{t=1}^T\eta_{J_t}^2,\label{eq:Q-budget}\\
 \mathcal D_{T,k}&=\E\sum_{t=1}^T
 \eta_k(H_t-4c_k)_+\ind\{J_t<k\}.\label{eq:D-budget}
\end{align}

\begin{theorem}[Thresholded shared-trajectory reduction]\label{thm:thresholded-reduction}
For every such selector and every $k$,
\begin{equation}\label{eq:thresholded-reduction}
 \frac1T\sum_{t=1}^T\E H_t
 \le4c_k+
 \frac{2\Delta+4\mathcal A_T+L\mathcal Q_T+\mathcal D_{T,k}}
 {\eta_kT}.
\end{equation}
In the normalized model, suppose a common selector satisfies
$\mathcal A_T+\mathcal D_{T,k}=O_p(1)$ for some grid index with
$s_k=a(p)+O_p(1/\log T)$.
Then a uniformly sampled iterate has expected gradient norm
$O_p(T^{-a(p)})$.
\end{theorem}
\begin{proof}
The normalized-direction inequality implies
\[
 \inner{h_t}{d_t}
 \ge\frac12H_t-2(e_t-H_t/4)_+
\]
whenever an update is made. Multiplying by $\eta_{J_t}$ also covers
the no-update action. Smoothness and telescoping therefore give
\begin{equation}\label{eq:weighted-descent}
 \E\sum_{t=1}^T\eta_{J_t}H_t
 \le2\Delta+4\mathcal A_T+L\mathcal Q_T.
\end{equation}
If $J_t\ge k$, monotonicity of the step sizes yields
$\eta_k(H_t-4c_k)_+\le\eta_{J_t}H_t$.
On the remaining rounds, this quantity is included in
$\mathcal D_{T,k}$. Consequently,
\[
 \eta_k\sum_{t=1}^T\E(H_t-4c_k)_+
 \le2\Delta+4\mathcal A_T+L\mathcal Q_T+\mathcal D_{T,k}.
\]
Use $H_t\le4c_k+(H_t-4c_k)_+$ and divide by $\eta_kT$.
For the final assertion,
$\mathcal Q_T\le T\eta_K^2\le T^{-1/2}$,
$\eta_kT=1/c_k$, and $c_k=\Theta_p(T^{-a(p)})$.
\end{proof}

The budget $\mathcal A_T$ charges direction errors exceeding a fraction
of the gradient norm. The quantity $\mathcal D_{T,k}$ charges rounds
where the gradient is above scale $c_k$ but the chosen step is smaller
than $\eta_k$. Bounding these two quantities is a sufficient condition;
no necessity statement is asserted.

\subsection{What observed variation can establish}

The variation statistic used by optimistic online learners such as
$\mathsf{OAdaGradR}$ \cite{liu2026oco} admits the following estimate.

\begin{lemma}[Variation moment]\label[lemma]{lem:variation}
Let $\cG$ be the history before a block. Query $g_0$ at a
$\cG$-measurable point $z_0$ and then query $g_i$ at predictable points
$z_i$, where $\norm{z_i-z_{i-1}}\le2D$ for $1\le i\le\tau$.
Assume \eqref{eq:oracle} at every query, including the initial query
conditionally on $\cG$. For
$V_\tau=\sum_{i=1}^{\tau}\norm{g_i-g_{i-1}}^2$, one has
\begin{equation}\label{eq:variation-moment}
 \E[V_\tau^{p/2}\mid\cG]\le C_p\tau(\sigma^p+(LD)^p).
\end{equation}
If $D=(N\tau)^{-1/2}$ and $B_\tau=\kappa\tau^{3/2}/\sqrt N$,
where $N,\kappa>0$, then
\begin{equation}\label{eq:variation-tail}
 \Prb(\sqrt{V_\tau}>B_\tau\mid\cG)
 \le C_p\kappa^{-p}
 \left[\sigma^pN^{p/2}\tau^{1-3p/2}+L^p\tau^{1-2p}\right].
\end{equation}
\end{lemma}
\begin{proof}
Since $p/2\le1$,
$V_\tau^{p/2}\le\sum_{i=1}^{\tau}\norm{g_i-g_{i-1}}^p$.
Write $g_i=\nabla F(z_i)+\xi_i$ and use
\[
 \norm{g_i-g_{i-1}}^p
 \le3^{p-1}\left(\norm{\xi_i}^p+\norm{\xi_{i-1}}^p+(2LD)^p\right).
\]
Conditional expectation, the tower property, and summation prove
\eqref{eq:variation-moment}. Markov's inequality at threshold
$B_\tau^p$, followed by substitution of $D$, proves
\eqref{eq:variation-tail}.
\end{proof}

For $L=\sigma=1$ and $N\ge1$, put $\tau_p(N)=N^{p/(3p-2)}$.
At a probe length $\tau\ge2^j\tau_p(N)$, $j\ge0$,
\eqref{eq:variation-tail} bounds the conditional rejection probability
by $C_p\kappa^{-p}2^{-j(3p-2)/2}$.
Thus a sufficiently long fresh probe has a small rejection probability.
This statement does not certify the accuracy of a block that passes
the test. For general $L,\sigma$, both terms in
\eqref{eq:variation-tail} must be retained.

\subsection{A limitation of local agreement certificates}

\begin{lemma}[Rare-outlier indistinguishability]\label[lemma]{lem:rare-outlier}
Fix $p\in(1,2]$, $0<\delta\le1/2$, and
$\mu=\tfrac12\delta^{(p-1)/p}$. Consider
\[
 P_0:\ Y=-\mu\ \text{almost surely},\qquad
 P_1:\ Y=\begin{cases}
 -\mu,&\text{with probability }1-\delta,\\
 (1-\delta)\mu/\delta,&\text{with probability }\delta.
 \end{cases}
\]
Their means are $-\mu$ and zero, respectively. Both have centered
$p$th moment at most one, and, for $N$ independent observations,
\[
 \operatorname{TV}(P_0^{\otimes N},P_1^{\otimes N})
 =1-(1-\delta)^N\le N\delta.
\]
\end{lemma}
\begin{proof}
The assertions for $P_0$ and the mean of $P_1$ are immediate. Moreover,
\[
 \E_{P_1}|Y|^p
 =2^{-p}\left[(1-\delta)\delta^{p-1}+(1-\delta)^p\right]
 \le2^{1-p}\le1.
\]
The only transcript under $P_0^{\otimes N}$ is the all-$-\mu$
transcript, whose probability under $P_1^{\otimes N}$ is
$(1-\delta)^N$. This proves the equality; Bernoulli's inequality
gives the upper bound.
\end{proof}

Set $\delta=1/w$ and $N=w\ge2$. With probability
$(1-1/w)^w$, which is bounded away from zero, every observation under
$P_1$ equals $-\mu$. All window means then agree and all observed
variation vanishes, yet the empirical mean differs from the true mean
by $\tfrac12w^{-b(p)}$.
The implication is limited to deterministic certification from observed
agreement. The example is consistent with the usual statistical error
scale and leaves open control of the budgets in
\cref{thm:thresholded-reduction}.


\begin{thebibliography}{99}

\bibitem{arjevani2023lower}
Yossi Arjevani, Yair Carmon, John C. Duchi, Dylan J. Foster,
Nathan Srebro, and Blake Woodworth.
\newblock Lower bounds for non-convex stochastic optimization.
\newblock \emph{Mathematical Programming}, 199(1--2):165--214, 2023.
\newblock \url{https://arxiv.org/abs/1912.02365}.

\bibitem{cutkosky2023o2nc}
Ashok Cutkosky, Harsh Mehta, and Francesco Orabona.
\newblock Optimal stochastic non-smooth non-convex optimization through
online-to-non-convex conversion.
\newblock In \emph{Proceedings of the 40th International Conference on
Machine Learning}, volume 202 of PMLR, pages 6643--6670, 2023.
\newblock \url{https://proceedings.mlr.press/v202/cutkosky23a.html}.

\bibitem{hubler2025normalization}
Florian H\"ubler, Ilyas Fatkhullin, and Niao He.
\newblock From gradient clipping to normalization for heavy tailed SGD.
\newblock In \emph{Proceedings of the 28th International Conference on
Artificial Intelligence and Statistics}, volume 258 of PMLR,
pages 2413--2421, 2025.
\newblock \url{https://proceedings.mlr.press/v258/hubler25a.html}.
\newblock Extended version: \url{https://arxiv.org/abs/2410.13849v3}.

\bibitem{khaled2024tuning}
Ahmed Khaled and Chi Jin.
\newblock Tuning-free stochastic optimization.
\newblock In \emph{Proceedings of the 41st International Conference on
Machine Learning}, volume 235 of PMLR, pages 23622--23661, 2024.
\newblock \url{https://proceedings.mlr.press/v235/khaled24a.html}.

\bibitem{liu2026oco}
Zijian Liu.
\newblock Online convex optimization with heavy tails: Old algorithms,
new regrets, and applications.
\newblock Extended version, arXiv:2508.07473v3, March 19, 2026.
\newblock \url{https://arxiv.org/abs/2508.07473v3}.
\newblock A shorter version appeared at ALT 2026. The theorem numbers
cited in this paper refer to the extended version.

\bibitem{liuzhou2025}
Zijian Liu and Zhengyuan Zhou.
\newblock Nonconvex stochastic optimization under heavy-tailed noises:
Optimal convergence without gradient clipping.
\newblock In \emph{International Conference on Learning Representations}, 2025.
\newblock \url{https://arxiv.org/abs/2412.19529v4}.

\bibitem{nguyen2023clipped}
Ta Duy Nguyen, Thien Hang Nguyen, Alina Ene, and Huy Le Nguyen.
\newblock Improved convergence in high probability of clipped gradient
methods with heavy tailed noise.
\newblock In \emph{Advances in Neural Information Processing Systems 36},
pages 24191--24222, 2023.

\bibitem{longtails2026stationarity}
Alexander Tyurin, QuocViet Pham, Valentin Khrulkov, and Ivan Oseledets.
\newblock LongTails: Automated research in optimization.
\newblock Online result: \emph{Parameter-Free Heavy-Tailed Stationarity},
2026. Consulted November 23, 2026.
\newblock \url{https://k3nfalt.github.io/longtails/setups/ParameterFreeHeavyTailedSmooth/}.
\newblock Algorithm and theorem sources:
\href{https://github.com/k3nfalt/longtails/blob/main/LongTails/Algorithms/ParameterFreeNsgdm.lean}{\texttt{ParameterFreeNsgdm.lean}} and
\href{https://github.com/k3nfalt/longtails/blob/main/LongTails/Theorems/ParameterFreeHeavyTailedNsgdm.lean}{\texttt{ParameterFreeHeavyTailedNsgdm.lean}}.

\bibitem{zhang2020adaptive}
Jingzhao Zhang, Sai Praneeth Karimireddy, Andreas Veit, Seungyeon Kim,
Sashank Reddi, Sanjiv Kumar, and Suvrit Sra.
\newblock Why are adaptive methods good for attention models?
\newblock In \emph{Advances in Neural Information Processing Systems 33},
pages 15383--15393, 2020.

\end{thebibliography}
\end{document}
